% !TeX root = ../组合数学笔记.tex
% 由主文档通过 \input 载入；请编译主文档。
\chapter{绪言与课程路线}
\label{ch:course-outline}

\section{研究对象与学习目标}

\emph{组合数学（combinatorics）}研究离散对象的计数、存在、构造及其结构。安排座位、分配物品、选择代表、统计路径、设计连接网络和识别对称图案，看似不同，却都需要先明确对象与限制，再寻找能够解释答案的方法。本书面向本科生，以有限集合、初等代数和基本证明方法为起点，沿着“对象建模---精确计数---结构分析---存在与构造”的路线组织内容，不把组合数学处理为孤立公式和解题技巧的汇编。

全书的学习目标是：能够把文字问题转化为无歧义的数学对象；根据约束选择分类、分步、双射、递推、生成函数或群作用等方法；独立完成主要结论的证明；区分精确公式、上下界、存在性证明与可执行算法。矩阵树定理和线性码专题需要线性代数，概率方法需要有限概率空间与期望，相关预备知识在附录中补充，不作为阅读开篇内容的障碍。

本书以定义、定理、证明和例题组成连续正文，围绕有限对象建立可核查的论证。核心结果给出证明，较长推导尽量由已学方法分解；所依赖的外部预备事实在使用处说明，不以证明思路冒充完整证明。章末习题用于检验模型、训练论证，附录给出代表性习题的解答。小规模计算能够核对例子，不能替代对任意参数成立的数学证明。

\section{全书大纲}

全书内容按知识依赖排列，各章在课程中的作用如下。
\begin{enumerate}
  \item 第\ref{ch:objects}章，绪论、组合对象与证明方法：集合语言、对象编码、基本计数原理、双射、归纳与构造。
  \item 第\ref{ch:elementary-counting}章，排列、组合与分配模型：有序与无序、允许重复与不允许重复、多重集合、圆排列和隔板法。
  \item 第\ref{ch:identities}章，二项式系数与组合恒等式：二项式与多项式定理、双重计数、带标记对象和格路径。
  \item 第\ref{ch:pigeonhole}章，抽屉原理与存在性证明：抽屉的选择、加强形式、余数模型、单调子序列和不变量。
  \item 第\ref{ch:inclusion-exclusion}章，容斥原理与二项式反演：交叠条件、错排、满射、有限上界分配与反演思想。
  \item 第\ref{ch:recurrences}章，递推关系与递归结构：组合分解、初值、线性递推、特征根和有限状态模型。
  \item 第\ref{ch:ogf}章，普通生成函数：形式幂级数、卷积、递推求解、受限分配和多变量标记。
  \item 第\ref{ch:egf}章，指数生成函数与集合分拆：标号结构、Stirling 数、Bell 数、置换循环和指数公式。
  \item 第\ref{ch:classical-structures}章，Catalan 数与整数分拆：括号、树与路径的对应，反射法、Ferrers 图及分拆恒等式。
  \item 第\ref{ch:graphs-trees}章，图、路径与树的计数：图的语言、连通性、Euler 道、树、Pr\"ufer 编码与生成树。
  \item 第\ref{ch:matching-coloring}章，匹配、染色与图的结构：Hall 定理、增广路、K\H{o}nig 定理、染色多项式和平面图入门。
  \item 第\ref{ch:symmetry}章，群作用与对称计数：轨道、稳定子、Burnside 引理、循环指标和 P\'olya 计数。
  \item 第\ref{ch:posets}章，偏序集与 M\"obius 反演：链与反链、关联代数、一般反演、Sperner 定理及格的入门。
  \item 第\ref{ch:extremal-ramsey}章，极值组合与 Ramsey 理论：极值函数、Mantel 定理、Tur\'an 图、Ramsey 数及上下界。
  \item 第\ref{ch:probabilistic}章，概率方法：随机结构、第一矩方法、随机二染色、删改法和条件期望法。
  \item 第\ref{ch:designs-codes}章，组合设计与纠错编码：区组设计、有限几何小例子、正交拉丁方、距离与编码界。
  \item 第\ref{ch:algorithms}章，组合算法与综合问题：枚举、动态规划、状态压缩、回溯，以及多种方法的比较与核验。
\end{enumerate}

\section{课程层次与知识依赖}

第\ref{ch:objects}至\ref{ch:inclusion-exclusion}章建立计数与证明基础；第\ref{ch:recurrences}至\ref{ch:classical-structures}章形成递推、生成函数和经典对象之间的联系；第\ref{ch:graphs-trees}至\ref{ch:symmetry}章转向离散结构与对称。以上十二章构成本书的主体，但矩阵树定理的完整证明、平面图的较深结论和完整 P\'olya 理论可以在首次学习时选读。第\ref{ch:posets}至\ref{ch:algorithms}章是五个拓展专题，不要求在一门短学期课程中全部讲完。

阅读时应保留以下依赖：普通生成函数依赖递推与基本计数，指数生成函数进一步依赖标号对象的分配；Catalan 数的闭式推导使用格路径与生成函数；匹配和极值图论需要图的基本语言；对称计数在章内补充有限群作用，不预设抽象代数课程；偏序反演从容斥和二项式反演出发；概率方法在使用 Ramsey 问题作为例子时衔接前一章。算法章中的基础枚举与动态规划可以提前配合相应计数章节阅读，设计与编码章则可独立选学。

基础学习路线以“每种方法能独立解决一个新问题”为验收标准，不以记住公式数量为目标。各章从具体问题引出概念和核心定理，再用直接应用、变式和反例辨明条件；章末练习覆盖概念辨析、计算建模、证明训练与综合拓展。涉及高级结果时，明确标注完整证明、证明思路或引用不证，不把未说明的结论当作前提。

\section{符号约定与参考范围}

记 $\mathbb Q,\mathbb R,\mathbb C$ 分别为有理数、实数、复数集，$\mathbb Z_{\ge0}$、$\mathbb Z_{>0}$ 分别为非负整数集和正整数集；对 $n\in\mathbb Z_{\ge0}$，记 $[n]=\{1,\ldots,n\}$，其中 $[0]=\varnothing$。有限集合 $A$ 的基数记为 $|A|$，其幂集记为 $2^A$。约定 $0!=1$，空和为 $0$，空积为 $1$；当 $n\ge0$ 而整数 $k$ 不满足 $0\le k\le n$ 时，置 $\binom nk=0$。计数问题中的参数范围由具体命题说明，不能依赖字母形状猜测。

本书区分\emph{标号（labeled）}与\emph{无标号（unlabeled）}对象，并为每个商计数问题明确等价关系。后文的 $[x^n]A(x)$ 表示形式幂级数 $A(x)$ 中 $x^n$ 的系数；图通常是有限简单无向图，需要允许重边、环或方向时另行声明。概率专题使用大写字母表示随机变量，小写字母表示其取值或固定参数。

课程范围参考 MIT 的本科组合分析课程资料\cite{mit-combinatorics}，问题组织可参考 Bogart 的开放教材\cite{bogart}，偏序、标号结构与群作用的拓展次序可与 Dartmouth 的代数组合课程资料\cite{dartmouth-combinatorics}对照。这些来源用于核对课程覆盖与提供后续阅读入口；本书的章节划分和教学安排另行组织，不复制来源中的成套习题。

